\documentclass[12pt,aspectratio=1610,professionalfont,serif,t]{beamer}
\usepackage[osf,p]{newpxtext}
\usepackage{newpxmath}
\usepackage{fontawesome5}

\usepackage{fancyvrb}
\fvset{%frame=single,
commandchars=\\\{\},
%framesep=1mm,
fontfamily=helvetica,
fontsize=\normalsize
}
%\usepackage[bitstream-charter]{mathdesign}
\usepackage{listings}
\lstset{commentstyle=\color{orange}, keywordstyle=\color{yellow} }


\usetheme{Boadilla}
\useoutertheme{split}
\usecolortheme{albatross}


\setbeamertemplate{blocks}[rounded][shadow=false]
\setbeamersize{text margin left=\dimexpr 1cm+1ex\relax, text margin right=\dimexpr 1cm+1ex\relax}
\addtobeamertemplate{headline}{}{\usebeamerfont{frametitle}\vspace*{1ex}}
\setbeamertemplate{navigation symbols}{}
\setbeamertemplate{page number in head/foot}{}
%\beamerdefaultoverlayspecification{<+->}


\setbeamercolor*{structure}{fg=green!60!black,bg=black}
\setbeamercolor*{normal text}{fg=green!60!black,bg=black}
\setbeamercolor{palette primary}{use={structure,normal text},fg=green,bg=structure.bg}
\setbeamercolor{palette secondary}{use={structure,normal text},fg=structure.fg,bg=structure.bg}
\setbeamercolor{palette tertiary}{use={structure,normal text},fg=structure.fg,bg=structure.bg}
\setbeamercolor{palette quaternary}{use={structure,normal text},fg=white,bg=green!30!black}
\setbeamercolor{palette sidebar quaternary}{use={structure,normal text},fg=white,bg=green!30!black}
\setbeamercolor{title in head/foot}{use={structure,normal text},fg=white,bg=green!30!black}

\setbeamercolor*{example text}{use={structure},fg=green!60!black}
\setbeamercolor*{block body}{use={structure},bg=structure.bg}
\setbeamercolor*{block body alerted}{use={structure},bg=structure.bg}
\setbeamercolor*{block body example}{use={structure},bg=structure.bg}
\setbeamercolor*{block title}{parent=structure,bg=structure.bg}
\setbeamercolor*{block title alerted}{use={structure,alerted text},fg=alerted text.fg,bg=structure.bg}
\setbeamertemplate{navigation symbols}{}
\setbeamertemplate{items}[square]
\setbeamercolor{item projected}{fg=white}
\setbeamercolor*{normal text}{use={structure},fg=white!90!blue,bg=structure.bg}
\setbeamercolor*{separation line}{}
\setbeamercolor*{fine separation line}{}
\setbeamercolor{alerted text}{use={structure},fg=structure.fg}

\usepackage[italian]{babel}
\usepackage[utf8]{inputenc}
\usepackage{pgf}
\usepackage{verbatim}
\usepackage{inconsolata}
\usepackage{listings}
\lstset{language=C, frame=single,
	basicstyle=\ttfamily,
	numbers=left, numberstyle=\tiny\color{gray},
	numbersep=5pt, fancyvrb=true
}

\usepackage[noend]{algorithm2e}
\usepackage{pgf}
\usepackage{tikz}
\usepackage{graphicx}
\usetikzlibrary{%
	arrows,
	arrows.meta,
	positioning,
	calc,
	backgrounds,
	chains,
	matrix,
	patterns,
	automata,
	fit,
	graphs,
	decorations,
	decorations.pathmorphing,
	decorations.pathreplacing,
	decorations.markings,
}

\usepgflibrary{shapes,shapes.geometric}


\usepackage{url}


\usepackage{xmpmulti}
\usepackage{transparent}
\graphicspath{{img/}}

\newcommand{\zero}{\ensuremath{\mathtt{0}}\xspace}
\newcommand{\uno}{\ensuremath{\mathtt{1}}\xspace}
\newcommand{\due}{\ensuremath{\mathtt{2}}\xspace}
\newcommand{\tre}{\ensuremath{\mathtt{3}}\xspace}
\newcommand{\kleene}{\ensuremath{{}^{\ast}}\xspace}
\xspaceaddexceptions{\zero \uno \due \tre \kleene}


\newcommand{\simplebox}[1]{
      \begin{tikzpicture}
              \node[text centered, rectangle, draw=black, thick, text width=\linewidth,align=center, rounded corners, minimum height=1em] (C) {#1};
      \end{tikzpicture}
}
%
% In questo file si mettono i comandi utili solo per questo insegnamento.
%
\author{Gianluca Della Vedova}
\title{Linguaggi e Computabilità}
\institute{Univ. Milano--Bicocca\\
	\texttt{https://www.unimib.it/gianluca-della-vedova}}
\date{\today}



\begin{document}


\begin{frame}[fragile]
	\frametitle{Linguaggi e Computabilità: Lezione 13}
\begin{enumerate}
\item
Passaggio Espressioni regolari a $\epsilon$-NFA
\end{enumerate}
\end{frame}


\begin{frame}
	\frametitle{Esercizio 1: minimizzare il DFA}
	\begin{center}
		\begin{tikzpicture}[enfa]
			\node[state, initial]   (q0) at (0,0)  {$q_0$};
			\node[state]            (q1) at (3,0)  {$q_1$};
			\node[state, accepting] (q2) at (6,0)  {$q_2$};
			\node[state, accepting] (q3) at (0,-3) {$q_3$};
			\node[state]            (q4) at (3,-3) {$q_4$};
			\node[state]            (q5) at (6,-3) {$q_5$};
			\path (q0) edge [loop above]        node       {\uno}  ()
			           edge                     node       {\zero} (q1)
			      (q1) edge [bend left=20]      node       {\zero} (q4)
			           edge [bend left=20]      node       {\uno}  (q2)
			      (q2) edge [bend left=20]      node       {\zero} (q1)
			           edge                     node       {\uno}  (q5)
			      (q3) edge [bend left=20]      node       {\zero} (q4)
			           edge                     node       {\uno}  (q0)
			      (q4) edge [bend left=20]      node       {\zero} (q1)
			           edge [bend left=20]      node       {\uno}  (q3)
			      (q5) edge                     node       {\zero} (q4)
			           edge [loop below]        node       {\uno}  ();
		\end{tikzpicture}
	\end{center}
% Soluzione.
% Tutti gli stati sono raggiungibili.
% Partizione iniziale: {q2, q3} (finali), {q0, q1, q4, q5}.
% Con 1: q1 -> q2 e q4 -> q3 (finali), q0 -> q0 e q5 -> q5 (non finali).
% Partizione: {q0, q5}, {q1, q4}, {q2, q3}, che è stabile.
% DFA minimo (3 stati): A = {q0, q5} iniziale, B = {q1, q4}, C = {q2, q3} finale.
% A: 0 -> B, 1 -> A;  B: 0 -> B, 1 -> C;  C: 0 -> B, 1 -> A.
% Riconosce le parole che terminano con 01.
\end{frame}


\begin{frame}
	\frametitle{Esercizio 2: minimizzare il DFA}
	\begin{center}
		\begin{tikzpicture}[enfa]
			\node[state, initial, accepting] (q0) at (0,0)  {$q_0$};
			\node[state]            (q1) at (3,0)    {$q_1$};
			\node[state, accepting] (q4) at (6,0)    {$q_4$};
			\node[state]            (q2) at (0,-3)   {$q_2$};
			\node[state]            (q3) at (3,-3)   {$q_3$};
			\node[state]            (q5) at (6,-3)   {$q_5$};
			\node[state]            (q6) at (9,-1.5) {$q_6$};
			\path (q0) edge                     node       {\zero} (q1)
			           edge [bend left=20]      node       {\uno}  (q2)
			      (q1) edge [bend left=20]      node       {\zero} (q4)
			           edge [bend left=20]      node       {\uno}  (q3)
			      (q4) edge [bend left=20]      node       {\zero} (q1)
			           edge [bend left=20]      node       {\uno}  (q5)
			      (q2) edge                     node[swap] {\zero} (q3)
			           edge [bend left=20]      node       {\uno}  (q0)
			      (q3) edge [bend left=20]      node       {\zero} (q5)
			           edge [bend left=20]      node       {\uno}  (q1)
			      (q5) edge [bend left=20]      node       {\zero} (q3)
			           edge [bend left=20]      node       {\uno}  (q4)
			      (q6) edge                     node[swap] {\zero} (q4)
			           edge                     node       {\uno}  (q5);
		\end{tikzpicture}
	\end{center}
% Soluzione.
% q6 non è raggiungibile da q0: si elimina.
% Partizione iniziale: {q0, q4} (finali), {q1, q2, q3, q5}.
% Con 0: q1 -> q4 (finale), q2 -> q3, q3 -> q5, q5 -> q3 (non finali).
% Con 1: q2 -> q0, q5 -> q4 (finali), q1 -> q3, q3 -> q1 (non finali).
% Partizione: {q0, q4}, {q1}, {q2, q5}, {q3}, che è stabile.
% DFA minimo (4 stati): A = {q0, q4} iniziale e finale, B = {q1}, C = {q2, q5}, D = {q3}.
% A: 0 -> B, 1 -> C;  B: 0 -> A, 1 -> D;  C: 0 -> D, 1 -> A;  D: 0 -> C, 1 -> B.
% Riconosce le parole con un numero pari di 0 e un numero pari di 1.
\end{frame}


\begin{frame}
	\frametitle{Esercizio 3: minimizzare il DFA}
	\begin{center}
		\begin{tikzpicture}[enfa]
			\node[state, initial]   (q0) at (0,0)  {$q_0$};
			\node[state]            (q1) at (3,0)  {$q_1$};
			\node[state, accepting] (q2) at (6,0)  {$q_2$};
			\node[state, accepting] (q3) at (6,-3) {$q_3$};
			\node[state]            (q4) at (0,-3) {$q_4$};
			\node[state]            (q5) at (3,-3) {$q_5$};
			\path (q0) edge                     node       {\uno}  (q1)
			           edge                     node[swap] {\zero} (q4)
			      (q1) edge [loop above]        node       {\uno}  ()
			           edge [bend left=20]      node       {\zero} (q2)
			      (q2) edge [bend left=20]      node       {\uno}  (q1)
			           edge                     node       {\zero} (q3)
			      (q3) edge [loop right]        node       {\zero} ()
			           edge                     node       {\uno}  (q1)
			      (q4) edge [bend left=20]      node       {\zero} (q5)
			           edge [loop below]        node       {\uno}  ()
			      (q5) edge [bend left=20]      node       {\zero} (q4)
			           edge [loop below]        node       {\uno}  ();
		\end{tikzpicture}
	\end{center}
% Soluzione.
% Tutti gli stati sono raggiungibili.
% Partizione iniziale: {q2, q3} (finali), {q0, q1, q4, q5}.
% Con 0: q1 -> q2 (finale), q0 -> q4, q4 -> q5, q5 -> q4 (non finali).
% Partizione: {q0, q4, q5}, {q1}, {q2, q3}.
% Con 1: q0 -> q1, mentre q4 -> q4 e q5 -> q5.
% Partizione: {q0}, {q1}, {q2, q3}, {q4, q5}, che è stabile.
% DFA minimo (4 stati): A = {q0} iniziale, B = {q1}, C = {q2, q3} finale,
% D = {q4, q5} (stato pozzo).
% A: 0 -> D, 1 -> B;  B: 0 -> C, 1 -> B;  C: 0 -> C, 1 -> B;  D: 0 -> D, 1 -> D.
% Riconosce le parole che iniziano con 1 e terminano con 0.
\end{frame}


\begin{frame}
	\frametitle{Esercizio 4: minimizzare il DFA}
	Stato iniziale $q_0$, stati finali $q_0, q_3, q_6$.
	\begin{center}
		\begin{tabular}{c|cc}
			         & \zero  & \uno  \\
			\hline
                        $\to *q_0$    & $q_3$  & $q_1$ \\
			$q_1$    & $q_4$  & $q_2$ \\
			$q_2$    & $q_5$  & $q_3$ \\
                        $*q_3$    & $q_0$  & $q_4$ \\
			$q_4$    & $q_1$  & $q_5$ \\
			$q_5$    & $q_2$  & $q_6$ \\
                        $*q_6$    & $q_6$  & $q_1$ \\
			$q_7$    & $q_0$  & $q_7$ \\
		\end{tabular}
	\end{center}
% Soluzione.
% q7 non è raggiungibile da q0: si elimina.
% Partizione iniziale: {q0, q3, q6} (finali), {q1, q2, q4, q5}.
% Con 1: q2 -> q3 e q5 -> q6 (finali), q1 -> q2 e q4 -> q5 (non finali).
% Partizione: {q0, q3, q6}, {q1, q4}, {q2, q5}, che è stabile.
% DFA minimo (3 stati): A = {q0, q3, q6} iniziale e finale, B = {q1, q4}, C = {q2, q5}.
% A: 0 -> A, 1 -> B;  B: 0 -> B, 1 -> C;  C: 0 -> C, 1 -> A.
% Riconosce le parole in cui il numero di 1 è multiplo di 3.
\end{frame}


\end{document}
